\documentclass[11pt,a4paper]{article}

% --- Essential Packages ---
\usepackage[utf8]{inputenc}
\usepackage[T1]{fontenc}
\usepackage[margin=1in]{geometry}
\usepackage{amsmath,amssymb,amsthm}
\usepackage{booktabs}
\usepackage{tabularx}
\usepackage{array}
\usepackage{graphicx}
\usepackage{xcolor}
\usepackage{fancyhdr}
\usepackage{listings}
\usepackage{enumitem}
\usepackage{tikz}
\usetikzlibrary{shapes.geometric, arrows.meta, positioning, calc}
\usepackage{microtype}
\usepackage{hyperref}

% --- Color Definitions ---
\definecolor{primaryblue}{RGB}{26, 86, 160}
\definecolor{darkslate}{RGB}{30, 41, 59}
\definecolor{lightgray}{RGB}{248, 250, 252}
\definecolor{bordergray}{RGB}{226, 232, 240}
\definecolor{codebg}{RGB}{245, 247, 250}
\definecolor{codecomment}{RGB}{100, 116, 139}
\definecolor{codekeyword}{RGB}{14, 116, 144}
\definecolor{codestring}{RGB}{180, 83, 9}

% --- Hyperref Setup ---
\hypersetup{
    colorlinks=true,
    linkcolor=primaryblue,
    citecolor=primaryblue,
    urlcolor=primaryblue,
    pdftitle={Algorithm Analysis, Visualization, and Empirical Benchmarking Platform},
    pdfauthor={Discrete Mathematics Academic Project},
    pdfsubject={DA-IICT Discrete Mathematics & Algorithmics}
}

% --- Header and Footer ---
\pagestyle{fancy}
\fancyhf{}
\renewcommand{\headrulewidth}{0.4pt}
\renewcommand{\footrulewidth}{0.4pt}
\fancyhead[L]{\small\textcolor{darkslate}{\textbf{DA-IICT} | Discrete Mathematics \& Algorithmics}}
\fancyhead[R]{\small\textcolor{darkslate}{Algorithm Analysis \& Benchmarking Platform}}
\fancyfoot[L]{\small\textcolor{darkslate}{Kenneth H. Rosen Curriculum Alignment}}
\fancyfoot[R]{\thepage}

% --- Theorem and Definition Environments ---
\theoremstyle{definition}
\newtheorem{definition}{Definition}[section]
\newtheorem{property}{Property}[section]

\theoremstyle{plain}
\newtheorem{theorem}{Theorem}[section]
\newtheorem{lemma}[theorem]{Lemma}
\newtheorem{corollary}[theorem]{Corollary}

\theoremstyle{remark}
\newtheorem{remark}{Remark}[section]

% --- Listings Setup for Python ---
\lstset{
    backgroundcolor=\color{codebg},
    basicstyle=\ttfamily\footnotesize\color{darkslate},
    keywordstyle=\color{codekeyword}\bfseries,
    commentstyle=\color{codecomment}\itshape,
    stringstyle=\color{codestring},
    numbers=left,
    numberstyle=\tiny\color{codecomment},
    stepnumber=1,
    numbersep=8pt,
    frame=single,
    frameround=tttt,
    rulecolor=\color{bordergray},
    breaklines=true,
    breakatwhitespace=true,
    tabsize=4,
    showstringspaces=false,
    captionpos=b
}

% --- Custom Currency Macros ---
\newcommand{\INR}{\text{INR~}}

% -------------------------------------------------------------
% DOCUMENT METADATA
% -------------------------------------------------------------
\title{
    \vspace{-1.5cm}
    \textbf{\Large Dhirubhai Ambani Institute of Information and Communication Technology}\\
    \vspace{0.3cm}
    \textbf{\LARGE Algorithm Analysis, Visualization, and Empirical Benchmarking Platform}\\
    \vspace{0.2cm}
    \large An Interactive Computational Engine for Discrete Mathematics \& Algorithmic Invariants
}

\author{
    \textbf{Name}: Aarsh Patel\\
    \textbf{Student ID}: 202612062\\
    \textbf{Course}: Discrete Mathematics (Algorithmics Module) \\
    \textbf{Curriculum Framework}: Kenneth H. Rosen, \textit{Discrete Mathematics and Its Applications} (8th Ed.) \\
    \textbf{Core Stack}: Python 3.13, Streamlit, Plotly Engine, NumPy, Pandas, Pytest
}

\date{Academic Year 2026--2027}

% -------------------------------------------------------------
\begin{document}
% -------------------------------------------------------------

\maketitle
\thispagestyle{fancy}

\begin{abstract}
The \textbf{Algorithm Analysis, Visualization, and Empirical Benchmarking Platform} is an interactive computational and pedagogical system engineered to bridge the analytical gap between theoretical algorithm complexity (asymptotic bounds, inductive loop invariants, structural correctness proofs) and empirical execution performance (wall-clock latency, memory shift overheads, comparison counts, and sensitivity to input distributions). Developed for the Discrete Mathematics and Algorithmics curriculum at DA-IICT, the platform features a dual-engine architecture: a high-precision, uninstrumented benchmarking core capable of scaling to arrays of $n = 50{,}000$ under microsecond-resolution hardware timers, alongside a deterministic step-by-step state snapshot generator for granular visualization. Covering 10 classical algorithms across Searching, Sorting, String Matching, Greedy Optimization, and Supplementary Discrete Mathematics, this work validates theoretical asymptotic predictions against empirical execution metrics over five distinct input distribution profiles, verified by an automated 18-case Pytest validation suite.
\end{abstract}

\vspace{0.5cm}
\tableofcontents
\newpage

% -------------------------------------------------------------
\section{Project Overview \& Problem Formulation}
% -------------------------------------------------------------

\subsection{Motivation and Problem Statement}
In introductory discrete mathematics and computer science curricula, algorithm analysis is predominantly taught as an abstract, mathematical discipline centered on asymptotic orders of growth ($O, \Omega, \Theta$). While this mathematical foundation is indispensable, undergraduate students frequently experience substantial cognitive disconnect when attempting to reconcile:
\begin{enumerate}[label=(\roman*)]
    \item \textbf{Asymptotic Bounds vs. Real Hardware Timers}: Theoretical notations suppress hardware constants and low-order terms, leading students to assume that an $O(n \log n)$ comparison algorithm must strictly outpace an $O(n^2)$ algorithm for all practical input scales, ignoring memory movement and cache locality.
    \item \textbf{Internal Invariant Dynamics}: Standard textbook descriptions provide static pseudocode. Students cannot directly visualize how loop invariants are maintained or violated during execution, nor how pointers partition the problem space.
    \item \textbf{Input Distribution Sensitivity}: The significant disparity between best-case ($\Omega$), worst-case ($O$), and average-case performance across structured inputs (such as nearly sorted datasets or keys with massive duplication) is often lost without interactive empirical exploration.
\end{enumerate}

\subsection{Core Objectives}
To address these analytical challenges, the platform was engineered to deliver:
\begin{itemize}
    \item \textbf{Dual-Engine Implementation}: Strict separation of algorithm logic into pure, zero-overhead benchmarking routines and state-capturing step generators.
    \item \textbf{Rigorously Correct Algorithms}: Implementations adhering strictly to the formal pseudocode specifications in Kenneth H. Rosen's \textit{Discrete Mathematics and Its Applications}.
    \item \textbf{Granular Step Tracing}: Capture of complete execution snapshots $\mathcal{S}_t$ at every atomic decision point, enabling forward/backward step-through, pointer visualizers, interval halving highlights, and scheduling timelines.
    \item \textbf{Controlled Empirical Profiling}: Reproducible benchmarks across customizable input distributions with multi-trial statistical averaging and hardware-invariant operation counters (comparisons, swaps, shifts).
\end{itemize}

% -------------------------------------------------------------
\section{Theoretical Foundations \& Mathematical Framework}
% -------------------------------------------------------------

\subsection{Rosen's Seven Properties of Algorithms}
According to Rosen (\S3.1), an algorithmic specification must satisfy seven fundamental properties. The platform explicitly demonstrates these properties across every module:
\begin{enumerate}
    \item \textbf{Input}: Specified instances from a well-defined domain (e.g., sequence $\langle a_1, \dots, a_n \rangle$, target element $x$, pattern string $P$).
    \item \textbf{Output}: Correctly mapped solution from the specified range (e.g., target index $i \in \{0, \dots, n-1\} \cup \{-1\}$, sorted permutation $A'$, minimal coin multiset, stable pairing $M$).
    \item \textbf{Definiteness}: Every computational step is strictly unambiguous.
    \item \textbf{Correctness}: Produces the mathematically correct output for every valid input instance.
    \item \textbf{Finiteness}: Guaranteed termination after a finite number of operations for any input.
    \item \textbf{Effectiveness}: Every instruction can be performed exactly in a finite duration by an abstract computing agent.
    \item \textbf{Generality}: The procedure is parameterized to solve all instances of the intended problem space.
\end{enumerate}

\subsection{Loop Invariants and Proof by Mathematical Induction}
The primary mathematical apparatus for proving the partial correctness of iterative algorithms is the \textbf{loop invariant}. A loop invariant is a formal predicate $P$ over program variables satisfying three conditions analogous to the principle of mathematical induction:
\begin{enumerate}
    \item \textbf{Initialization (Base Case)}: $P$ holds true prior to the first iteration of the loop.
    \item \textbf{Maintenance (Inductive Step)}: If $P$ is true prior to iteration $k$, it remains true prior to iteration $k+1$.
    \item \textbf{Termination}: Upon termination, the invariant provides a guarantee establishing the algorithm's correctness.
\end{enumerate}

\begin{theorem}[Correctness of Insertion Sort]
Insertion Sort correctly sorts any array $A[0 \dots n-1]$ of $n$ elements.
\end{theorem}

\begin{proof}
Let the loop invariant predicate $P(i)$ be: \textit{At the start of iteration $i$ of the outer loop ($1 \le i \le n$), the subarray $A[0 \dots i-1]$ consists of the elements originally in $A[0 \dots i-1]$, but in sorted non-decreasing order ($A[0] \le A[1] \le \dots \le A[i-1]$).}

\begin{itemize}
    \item \textbf{Initialization}: Prior to the first iteration ($i = 1$), the subarray is $A[0 \dots 0]$. A singleton array is trivially sorted. Thus, $P(1)$ holds.
    \item \textbf{Maintenance}: Suppose $P(i)$ holds at the start of iteration $i$. The algorithm selects $\text{key} = A[i]$ and compares it right-to-left with elements $A[j]$ for $j = i-1$ down to $0$. While $A[j] > \text{key}$, $A[j]$ is shifted to $A[j+1]$. This preserves the relative sorted order of all elements greater than $\text{key}$. When the condition fails or $j < 0$, $\text{key}$ is placed into $A[j+1]$. The new subarray $A[0 \dots i]$ contains the elements of $A[0 \dots i-1]$ plus $\text{key}$ in non-decreasing order. Incrementing $i$ to $i+1$ preserves $P(i+1)$.
    \item \textbf{Termination}: The outer loop terminates when $i = n$. By substituting $i = n$ into $P(i)$, the subarray $A[0 \dots n-1]$ consists of all original elements in sorted order. Hence, the algorithm is correct.
\end{itemize}
\end{proof}

\subsection{Formal Asymptotic Notations}
The platform incorporates formal definitions from Rosen (\S3.2) to evaluate resource growth rates:

\begin{definition}[Big-$O$ Notation]
Let $f(n)$ and $g(n)$ be functions from $\mathbb{Z}^+$ or $\mathbb{R}^+$ to $\mathbb{R}$. We write $f(n) \in O(g(n))$ if there exist positive real constants $c$ and $k$ such that:
\begin{equation}
|f(n)| \le c \cdot |g(n)| \quad \forall n > k
\end{equation}
\end{definition}

\begin{definition}[Big-$\Omega$ Notation]
We write $f(n) \in \Omega(g(n))$ if there exist positive real constants $c$ and $k$ such that:
\begin{equation}
|f(n)| \ge c \cdot |g(n)| \quad \forall n > k
\end{equation}
\end{definition}

\begin{definition}[Big-$\Theta$ Notation]
We write $f(n) \in \Theta(g(n))$ if $f(n) \in O(g(n))$ and $f(n) \in \Omega(g(n))$, indicating the existence of positive constants $c_1, c_2,$ and $k$ such that:
\begin{equation}
c_1 \cdot |g(n)| \le |f(n)| \le c_2 \cdot |g(n)| \quad \forall n > k
\end{equation}
\end{definition}

\begin{definition}[Little-$o$ Notation]
We write $f(n) \in o(g(n))$ if:
\begin{equation}
\lim_{n \to \infty} \frac{f(n)}{g(n)} = 0
\end{equation}
\end{definition}

\subsection{Greedy Paradigms: Choice Property and Optimal Substructure}

\subsubsection{Cashier's Change-Making Algorithm}
Given a target value $V \in \mathbb{Z}^+$ and sorted coin denominations $c_1 > c_2 > \dots > c_k = 1$, the greedy choice picks the largest denomination $c_i \le V$ and sets $V \leftarrow V - c_i$ iteratively.

\begin{property}[Canonical vs. Non-Canonical Systems]
In a canonical coin system (such as US coins $\{25, 10, 5, 1\}$ or Indian currency $\{500, 200, 100, 50, 20, 10, 5, 2, 1\}$), the greedy choice provably yields the global minimum number of coins.
\end{property}

\begin{remark}[Failure on Non-Canonical Systems]
Consider denominations $C = \{4, 3, 1\}$ and change $V = 6$:
\begin{itemize}
    \item \textbf{Greedy Choice}: Selects $4$, remaining $2$; then selects $1$ and $1 \implies \{4, 1, 1\}$ ($3$ coins).
    \item \textbf{Optimal Solution}: Selects $3$ and $3 \implies \{3, 3\}$ ($2$ coins).
\end{itemize}
The visualizer demonstrates this boundary condition clearly by allowing arbitrary user-defined denomination sets.
\end{remark}

\subsubsection{Activity Selection Problem (Interval Scheduling)}
Given $n$ activities with intervals $[s_i, f_i)$, the objective is to maximize the cardinality of a mutually compatible subset.

\begin{theorem}[Optimality of Earliest Finish Time First]
The greedy algorithm that repeatedly selects the activity with the minimal finish time compatible with previously scheduled activities yields an optimal schedule.
\end{theorem}

\begin{proof}
Let $A = \langle a_1, a_2, \dots, a_k \rangle$ be the schedule chosen by the greedy algorithm ordered by finish times, and let $O = \langle o_1, o_2, \dots, o_m \rangle$ be an optimal schedule. We prove by induction that $f(a_r) \le f(o_r)$ for all $1 \le r \le k$.
\begin{itemize}
    \item \textbf{Base Case ($r=1$)}: The greedy algorithm picks $a_1$ as the activity with the global earliest finish time. Hence, $f(a_1) \le f(o_1)$.
    \item \textbf{Inductive Step}: Assume $f(a_{r-1}) \le f(o_{r-1})$. Since $o_r$ is compatible with $o_{r-1}$, we have $s(o_r) \ge f(o_{r-1}) \ge f(a_{r-1})$. Thus, $o_r$ is compatible with $a_{r-1}$. Because the greedy algorithm selects $a_r$ as the earliest-finishing compatible activity, $f(a_r) \le f(o_r)$.
    \item \textbf{Optimality}: If $m > k$, then $o_{k+1}$ exists and is compatible with $o_k$. Since $f(a_k) \le f(o_k)$, $o_{k+1}$ must also be compatible with $a_k$. This contradicts the termination condition of the greedy algorithm. Therefore, $k = m$, and the greedy selection is optimal.
\end{itemize}
\end{proof}

% -------------------------------------------------------------
\section{System Architecture \& Software Engineering Design}
% -------------------------------------------------------------

\subsection{Architecture Overview}
The platform is organized into three strictly decoupled layers, as shown in Figure~\ref{fig:arch}.

\begin{figure}[htbp]
\centering
\begin{tikzpicture}[
    box/.style={rectangle, draw=primaryblue, fill=codebg, thick, minimum width=3.2cm, minimum height=0.9cm, align=center, rounded corners=2pt},
    layerbox/.style={rectangle, draw=darkslate!40, fill=lightgray, dashed, inner sep=0.35cm, rounded corners=4pt},
    arrow/.style={-{Stealth[length=2.5mm]}, thick, primaryblue}
]

% Presentation Layer
\node[layerbox] (presentation) at (0, 3.8) {
    \begin{tabular}{c}
        \textbf{\large Presentation Layer (Streamlit \& Plotly Engine)} \\
        \begin{tikzpicture}
            \node[box] (nav) at (-4, 0) {Sidebar View Router};
            \node[box] (viz) at (-0.5, 0) {Visualizers\\(Step-by-Step)};
            \node[box] (bench) at (3, 0) {Benchmarking\\Studio};
            \node[box] (growth) at (6.5, 0) {Complexity\\Curves};
        \end{tikzpicture}
    \end{tabular}
};

% Engine Layer
\node[layerbox] (engines) at (0, 1.2) {
    \begin{tabular}{c}
        \textbf{\large Computational Engines} \\
        \begin{tikzpicture}
            \node[box] (stepgen) at (-3.5, 0) {Step Generators\\$\mathcal{S}_t$ Snapshots};
            \node[box] (purealgo) at (0.5, 0) {Pure Benchmarking\\Cores};
            \node[box] (datagen) at (4.5, 0) {Data Generator\\(5 Distributions)};
        \end{tikzpicture}
    \end{tabular}
};

% Domain Layer
\node[layerbox] (domain) at (0, -1.4) {
    \begin{tabular}{c}
        \textbf{\large Algorithmic Core Modules} \\
        \begin{tikzpicture}
            \node[box] (search) at (-5, 0) {Searching\\(Linear, Binary)};
            \node[box] (sort) at (-1.7, 0) {Sorting\\(4 Algorithms)};
            \node[box] (match) at (1.5, 0) {String Matching\\(Naive Shift)};
            \node[box] (greedy) at (4.8, 0) {Greedy \& Discrete\\(Change, Gale-Shapley)};
        \end{tikzpicture}
    \end{tabular}
};

% Connecting Arrows
\draw[arrow] (viz.south) -- (stepgen.north);
\draw[arrow] (bench.south) -- (purealgo.north);
\draw[arrow] (bench.south) -- (datagen.north);
\draw[arrow] (stepgen.south) -- (sort.north);
\draw[arrow] (purealgo.south) -- (search.north);
\draw[arrow] (purealgo.south) -- (greedy.north);

\end{tikzpicture}
\caption{Three-Tier Architecture of the Platform.}
\label{fig:arch}
\end{figure}

\subsection{Separation of Concerns: The Dual-Engine Pattern}
Embedding step-recording and visualization hooks into algorithms severely distorts benchmarking accuracy due to Python's dictionary allocations and list cloning overheads. To prevent this, every algorithm is implemented as two separate functions:
\begin{enumerate}
    \item \textbf{Pure Engine} (\texttt{algo(arr, ...)}): Executes pure mathematical comparisons and assignments. Returns the result alongside exact hardware-invariant operation counters (comparisons, swaps).
    \item \textbf{Step Engine} (\texttt{algo\_steps(arr, ...)}): Executes the algorithm while yielding a comprehensive list of discrete state dictionaries $\mathcal{S}_t$ used exclusively by the Plotly visualization engine.
\end{enumerate}

\subsection{State Snapshot Model}
Each state snapshot $\mathcal{S}_t$ forms an immutable record defined as:
\begin{equation}
\mathcal{S}_t = \left\langle \text{array}_t, \mathcal{I}_{\text{active}}, \mathcal{I}_{\text{mutated}}, \mathcal{I}_{\text{sorted}}, \text{counters}_t, \text{description}_t \right\rangle
\end{equation}
This design guarantees reproducible step navigation, enabling forward and backward time-travel stepping without re-running the algorithm.

% -------------------------------------------------------------
\section{Algorithmic Suite: Design, Analysis \& Implementation}
% -------------------------------------------------------------

\subsection{Searching Algorithms}

\subsubsection{Linear Search}
Sequentially scans array $A[0 \dots n-1]$ until target $x$ is encountered or the end of the array is reached.
\begin{itemize}
    \item \textbf{Best-Case Complexity}: $O(1)$ comparisons ($x = A[0]$).
    \item \textbf{Average-Case Complexity}: $\frac{n+1}{2} \in O(n)$ comparisons under uniform target distribution.
    \item \textbf{Worst-Case Complexity}: $n \in O(n)$ comparisons ($x \notin A$ or $x = A[n-1]$).
    \item \textbf{Space Complexity}: $O(1)$ auxiliary.
\end{itemize}

\subsubsection{Binary Search}
Requires sorted input $A[0 \le \dots \le n-1]$. Halves the search range at each step:
\begin{equation}
T(n) = T\left(\left\lfloor \frac{n}{2} \right\rfloor\right) + O(1) \implies T(n) \in \Theta(\log n)
\end{equation}
\begin{itemize}
    \item \textbf{Best-Case Complexity}: $O(1)$ comparisons ($x = A[mid]$ on step 1).
    \item \textbf{Worst-Case Complexity}: $\lfloor \log_2 n \rfloor + 1 \in \Theta(\log n)$ comparisons.
    \item \textbf{Space Complexity}: $O(1)$ auxiliary (iterative formulation).
\end{itemize}

\subsection{Sorting Algorithms}

\subsubsection{Bubble Sort (with Early Termination)}
Compares adjacent pairs and bubbles the maximum unsorted value to the right boundary. The boolean flag \texttt{swapped} detects early sorted state.
\begin{itemize}
    \item \textbf{Best-Case}: $n - 1 \in \Theta(n)$ comparisons, $0$ swaps (pre-sorted array).
    \item \textbf{Worst-Case}: $\frac{n(n-1)}{2} \in \Theta(n^2)$ comparisons and swaps (reverse-sorted array).
    \item \textbf{Stability}: Stable (does not swap equal keys: strictly $A[j] > A[j+1]$).
\end{itemize}

\subsubsection{Insertion Sort}
Builds the sorted prefix $A[0 \dots i-1]$ incrementally by shifting elements greater than $\text{key} = A[i]$ rightward.
\begin{itemize}
    \item \textbf{Best-Case}: $\Theta(n)$ comparisons, $\Theta(1)$ shifts (pre-sorted array).
    \item \textbf{Worst-Case}: $\frac{n(n-1)}{2} \in \Theta(n^2)$ comparisons and shifts (reverse-sorted array).
    \item \textbf{Stability}: Stable.
\end{itemize}

\subsubsection{Selection Sort}
Finds the minimal element in $A[i \dots n-1]$ and swaps it into index $i$.
\begin{itemize}
    \item \textbf{Comparisons Invariant}: Exactly $\sum_{i=0}^{n-2} (n - 1 - i) = \frac{n(n-1)}{2} \in \Theta(n^2)$ comparisons across \textit{all} cases.
    \item \textbf{Swaps}: Strictly $\le n - 1 \in O(n)$ swaps.
    \item \textbf{Stability}: Unstable in standard swap form (e.g., $[2_a, 2_b, 1]$ swaps $2_a$ past $2_b$).
\end{itemize}

\subsubsection{Binary Insertion Sort}
Replaces the linear prefix scan with binary search to find the insertion index in $O(\log i)$ comparisons.
\begin{itemize}
    \item \textbf{Comparisons}: $\sum_{i=1}^{n-1} \lceil \log_2(i+1) \rceil \in \Theta(n \log n)$ comparisons for all cases.
    \item \textbf{Memory Shifts}: Requires $O(n^2)$ physical shifts in contiguous memory.
    \item \textbf{Total Worst-Case Time}: $\Theta(n^2)$.
    \item \textbf{Stability}: Stable (binary search targets rightmost insertion position).
\end{itemize}

\subsection{Naive String Matching}
Slides pattern $P[0 \dots m-1]$ across text $T[0 \dots n-1]$ over shifts $s \in [0 \dots n-m]$.
\begin{itemize}
    \item \textbf{Best-Case Time}: $O(n)$ when $P[0]$ does not appear in $T$.
    \item \textbf{Worst-Case Time}: $O((n - m + 1) \cdot m) \in O(n \cdot m)$ on pathological inputs (e.g., $T = \text{``AAAAAAAAAA''}, P = \text{``AAAB''}$).
    \item \textbf{Space}: $O(1)$ auxiliary.
\end{itemize}

\subsection{Supplementary Discrete Algorithms}

\subsubsection{Boyer-Moore Majority Vote Algorithm}
Finds the majority element ($> \lfloor n/2 \rfloor$ occurrences) in streaming $O(n)$ time and $O(1)$ auxiliary space via paired cancellation:
\begin{enumerate}
    \item \textbf{Phase 1 (Candidate Selection)}: Maintains a running counter. Counter increments on candidate matches and decrements on mismatches; when counter reaches zero, the current element becomes the new candidate.
    \item \textbf{Phase 2 (Verification)}: Scans the array once to verify if the frequency strictly exceeds $n/2$.
\end{enumerate}

\subsubsection{Gale-Shapley Stable Marriage Algorithm}
Computes a stable bipartite matching in $O(n^2)$ proposals. Guarantees:
\begin{itemize}
    \item \textbf{Termination}: Terminates in at most $n^2$ proposals.
    \item \textbf{Stability}: No blocking pairs exist upon termination.
    \item \textbf{Proposer-Optimality}: Every proposer is paired with their best valid partner in any stable matching.
\end{itemize}

% -------------------------------------------------------------
\section{Master Complexity \& Behavioral Matrix}
% -------------------------------------------------------------

Table~\ref{tab:master_complexity} summarizes the theoretical complexity, stability, and memory properties across all 10 algorithms implemented in the platform.

\begin{table}[htbp]
\centering
\small
\caption{Comprehensive Complexity and Property Matrix.}
\label{tab:master_complexity}
\begin{tabularx}{\textwidth}{l l c c c c c c}
\toprule
\textbf{Algorithm} & \textbf{Category} & \textbf{Best Time} & \textbf{Avg Time} & \textbf{Worst Time} & \textbf{Aux Space} & \textbf{Stable} & \textbf{In-Place} \\
\midrule
Linear Search & Searching & $O(1)$ & $O(n)$ & $O(n)$ & $O(1)$ & N/A & Yes \\
Binary Search & Searching & $O(1)$ & $O(\log n)$ & $O(\log n)$ & $O(1)$ & N/A & Yes \\
Bubble Sort & Sorting & $O(n)$ & $O(n^2)$ & $O(n^2)$ & $O(1)$ & Yes & Yes \\
Insertion Sort & Sorting & $O(n)$ & $O(n^2)$ & $O(n^2)$ & $O(1)$ & Yes & Yes \\
Selection Sort & Sorting & $O(n^2)$ & $O(n^2)$ & $O(n^2)$ & $O(1)$ & No & Yes \\
Binary Insertion Sort & Sorting & $O(n \log n)$ & $O(n^2)$ & $O(n^2)$ & $O(1)$ & Yes & Yes \\
Naive String Match & String & $O(n)$ & $O(n+m)$ & $O(n \cdot m)$ & $O(1)$ & N/A & Yes \\
Cashier's Change & Greedy & $O(k)$ & $O(k)$ & $O(k)$ & $O(1)$ & N/A & Yes \\
Activity Selection & Greedy & $O(n \log n)$ & $O(n \log n)$ & $O(n \log n)$ & $O(n)$ & N/A & No \\
Boyer-Moore Majority & Discrete & $O(n)$ & $O(n)$ & $O(n)$ & $O(1)$ & N/A & Yes \\
Gale-Shapley Matching & Discrete & $O(n)$ & $O(n \log n)$ & $O(n^2)$ & $O(n^2)$ & Yes & No \\
\bottomrule
\end{tabularx}
\end{table}

% -------------------------------------------------------------
\section{Empirical Benchmarking \& Experimental Evaluation}
% -------------------------------------------------------------

\subsection{Benchmarking Methodology}
To ensure rigorous empirical evaluation, the platform follows four scientific protocols:
\begin{enumerate}
    \item \textbf{Seed-Matched Datasets}: In any comparison run of size $n$, all evaluated algorithms receive exact deep copies of the identical pseudo-random array generated by a fixed NumPy seed.
    \item \textbf{Multi-Trial Averaging}: Each measurement is repeated over $k = 5$ independent trials to minimize OS thread scheduling and background task variance.
    \item \textbf{High-Resolution Timing}: Execution elapsed time is measured with Python's monotonic clock \texttt{time.perf\_counter()}, providing microsecond precision.
    \item \textbf{Hardware-Invariant Metrics}: In addition to clock time, atomic comparison and swap/shift operations are tracked independently.
\end{enumerate}

\subsection{Sorting Benchmarks Across Input Distributions}
Evaluation was performed across five synthetic input distributions ($n \in [50, 500]$):
\begin{itemize}
    \item \textbf{Random Uniform}: Independent and identically distributed uniform random integers.
    \item \textbf{Sorted}: Monotonically non-decreasing integers.
    \item \textbf{Reverse Sorted}: Pathological descending order.
    \item \textbf{Nearly Sorted}: Pre-sorted array with $8\%$ random index perturbations.
    \item \textbf{Many Duplicates}: Array sampled from only 5 unique integer values.
\end{itemize}

\subsubsection{Key Empirical Findings}
\begin{enumerate}
    \item \textbf{Nearly Sorted Efficiency}: Insertion Sort decisively outpaces Bubble and Selection Sort, executing close to its theoretical $O(n)$ best-case. Selection Sort shows zero improvement, executing the exact same $\frac{n(n-1)}{2}$ comparisons.
    \item \textbf{The Binary Insertion Sort Paradox}: On random arrays of size $n = 400$, Binary Insertion Sort slashes comparisons by \textbf{$\approx 80\%$} compared to standard Insertion Sort ($\approx 3{,}200$ comparisons vs. $\approx 40{,}000$). However, its wall-clock time is nearly identical to standard Insertion Sort. This empirically demonstrates that contiguous array pointer shifts ($O(n^2)$ in Python lists) dominate execution latency over comparison savings.
    \item \textbf{Selection Sort Swap Advantage}: On reverse-sorted inputs, Bubble and Insertion Sort make $\approx \frac{n(n-1)}{2}$ swaps/shifts. Selection Sort performs strictly $n - 1$ swaps, demonstrating why Selection Sort can outperform Bubble Sort in hardware clock time despite identical $O(n^2)$ comparison counts.
\end{enumerate}

\subsection{Searching Benchmarks: Linear vs. Binary Search}
Table~\ref{tab:search_bench} documents the scaling divergence between Linear Search ($O(n)$) and Binary Search ($O(\log n)$) up to $n = 50{,}000$.

\begin{table}[htbp]
\centering
\caption{Empirical Search Operation Counts vs. Theoretical Bounds.}
\label{tab:search_bench}
\begin{tabularx}{0.9\textwidth}{r c c c c c}
\toprule
\textbf{Array Size ($n$)} & \textbf{Linear (Avg)} & \textbf{Linear (Worst)} & \textbf{Binary (Avg)} & \textbf{Binary (Worst)} & \textbf{Theory $\lceil \log_2(n+1) \rceil$} \\
\midrule
100    & $\approx 50$    & 100    & $\approx 6$  & 7  & 7 \\
500    & $\approx 250$   & 500    & $\approx 8$  & 9  & 9 \\
1,000  & $\approx 500$   & 1,000  & $\approx 9$  & 10 & 10 \\
5,000  & $\approx 2,500$ & 5,000  & $\approx 11$ & 13 & 13 \\
10,000 & $\approx 5,000$ & 10,000 & $\approx 12$ & 14 & 14 \\
50,000 & $\approx 25,000$& 50,000 & $\approx 14$ & 16 & 16 \\
\bottomrule
\end{tabularx}
\end{table}

At $n = 50{,}000$, Linear Search requires an average of $25{,}000$ operations, whereas Binary Search never exceeds \textbf{16 comparisons}, executing in $< 0.005\text{ ms}$.

% -------------------------------------------------------------
\section{Verification \& Quality Assurance}
% -------------------------------------------------------------

The codebase incorporates an automated \textbf{Pytest} testing suite guaranteeing algorithmic correctness, state snapshot integrity, and metadata adherence.

\begin{lstlisting}[language=bash, caption={Pytest Execution Command and Output Summary}]
$env:PYTHONPATH="Dashboard"
python -m pytest Dashboard/tests -v

============================= test session starts =============================
platform win32 -- Python 3.13.17, pytest-9.2.0, pluggy-1.7.0
collected 18 items

Dashboard/tests/test_algorithms.py::test_linear_search PASSED            [  5%]
Dashboard/tests/test_algorithms.py::test_binary_search PASSED            [ 11%]
Dashboard/tests/test_algorithms.py::test_sorting_algorithms[bubble] PASSED [ 16%]
Dashboard/tests/test_algorithms.py::test_sorting_algorithms[insertion] PASSED [ 22%]
Dashboard/tests/test_algorithms.py::test_sorting_algorithms[selection] PASSED [ 27%]
Dashboard/tests/test_algorithms.py::test_sorting_algorithms[bin_insert] PASSED [ 33%]
Dashboard/tests/test_algorithms.py::test_sorting_steps PASSED            [ 38%]
Dashboard/tests/test_algorithms.py::test_naive_string_match PASSED       [ 44%]
Dashboard/tests/test_algorithms.py::test_cashiers_algorithm PASSED       [ 50%]
Dashboard/tests/test_algorithms.py::test_activity_selection PASSED       [ 55%]
Dashboard/tests/test_algorithms.py::test_boyer_moore_majority PASSED     [ 61%]
Dashboard/tests/test_algorithms.py::test_gale_shapley PASSED             [ 66%]
Dashboard/tests/test_algorithms.py::test_data_generator PASSED           [ 72%]
Dashboard/tests/test_analysis.py::test_complexity_metadata PASSED        [ 77%]
Dashboard/tests/test_analysis.py::test_asymptotic_notations PASSED       [ 83%]
Dashboard/tests/test_analysis.py::test_growth_data_calculation PASSED    [ 88%]
Dashboard/tests/test_analysis.py::test_sorting_benchmark PASSED          [ 94%]
Dashboard/tests/test_analysis.py::test_searching_benchmark PASSED        [100%]

============================= 18 passed in 0.92s ==============================
\end{lstlisting}

% -------------------------------------------------------------
\section{Educational Significance, Limitations \& Future Scope}
% -------------------------------------------------------------

\subsection{Educational Significance}
This platform transforms static discrete mathematics theorems into tangible, interactive laboratories. Students can directly verify:
\begin{itemize}
    \item How the inductive hypothesis of a loop invariant holds across each inner loop iteration.
    \item Why non-canonical currency denominations break greedy change-making.
    \item How theoretical bounds materialize on real computer architectures.
\end{itemize}

\subsection{Limitations}
\begin{itemize}
    \item \textbf{Snapshot Memory Footprint}: Generating step snapshots has space complexity $O(\text{steps} \cdot n)$. For this reason, the step visualizer is capped at $n \le 100$ in the UI, while pure benchmarking operates up to $n = 50{,}000$.
    \item \textbf{Interpreted Execution Overhead}: Python's bytecode interpretation introduces interpreter-level constant overheads absent in compiled C++ or Rust implementations.
\end{itemize}

\subsection{Future Roadmap}
\begin{itemize}
    \item \textbf{Divide-and-Conquer Sorting}: Integration of Merge Sort and Quick Sort (Lomuto and Hoare partitioning) with interactive recursion tree renderers.
    \item \textbf{Graph Algorithms (Rosen Chapter 10)}: Inclusion of Dijkstra's Shortest Path, Prim's, and Kruskal's Minimum Spanning Tree algorithms.
    \item \textbf{Dynamic Programming Comparison}: Side-by-side execution of 0/1 Knapsack (DP) versus Fractional Knapsack (Greedy) to visually illustrate optimal substructure failures.
\end{itemize}

% -------------------------------------------------------------
\section{Conclusion}
% -------------------------------------------------------------
The \textbf{Algorithm Analysis, Visualization, and Empirical Benchmarking Platform} successfully unifies theoretical discrete mathematics with empirical software engineering. By implementing pure benchmarking engines alongside discrete state snapshot generators, the platform provides students and researchers with a mathematically rigorous, interactive environment for exploring algorithm mechanics, loop invariants, asymptotic bounds, and real-world execution characteristics.

% -------------------------------------------------------------
% REFERENCES
% -------------------------------------------------------------
\begin{thebibliography}{9}

\bibitem{rosen2019}
K.~H. Rosen,
\textit{Discrete Mathematics and Its Applications},
8th~ed., McGraw-Hill Education, New York, 2019.

\bibitem{clrs2022}
T.~H. Cormen, C.~E. Leiserson, R.~L. Rivest, and C.~Stein,
\textit{Introduction to Algorithms},
4th~ed., MIT Press, Cambridge, MA, 2022.

\bibitem{boyer1991}
R.~S. Boyer and J.~S. Moore,
``MJRTY---A Fast Majority Vote Algorithm,''
in \textit{Automated Reasoning: Essays in Honor of Woody Bledsoe},
Springer, Dordrecht, 1991, pp.~105--117.

\bibitem{galeshapley1962}
D.~Gale and L.~S. Shapley,
``College Admissions and the Stability of Marriage,''
\textit{The American Mathematical Monthly},
vol.~69, no.~1, pp.~9--15, 1962.

\bibitem{streamlit2024}
Streamlit Inc.,
\textit{Streamlit Documentation and API Reference (v1.40+)},
2024. Available: \url{https://docs.streamlit.io}

\bibitem{plotly2024}
Plotly Technologies Inc.,
\textit{Plotly Python Graphing Library Reference (v5.24+)},
2024. Available: \url{https://plotly.com/python}

\bibitem{GitHub Repo}
Github Repo:
\hyperlink{https://github.com/aarshpatel1/DiscreteMathsAlgorithmProject}{https://github.com/aarshpatel1/DiscreteMathsAlgorithmProject}
\end{thebibliography}

\end{document}
